\reading{MR-6}{Validating Bank Holding Companies' Value-at-Risk Models for Market Risk}{defer}{0}
% ==========================================================

\begin{mrkeybox}[title={Learning objectives in this reading}]
\small
\begin{enumerate}[label=\textbf{\alph*.},leftmargin=1.7em]
\item Describe some important considerations for a bank in assessing the conceptual soundness of
      a VaR model during the validation process.
\item Explain how to conduct sensitivity analysis for a VaR model, and describe the potential
      benefits and challenges of performing such an analysis.
\item Describe the challenges a financial institution could face when calculating confidence
      intervals for VaR.
\item Discuss the challenges in benchmarking VaR models and various approaches proposed to
      overcome them.
\end{enumerate}
\end{mrkeybox}

\section{Basics}

Banks use Value at Risk (VaR) models to measure the potential loss in a portfolio over a
specified time period with a given probability. These models play a crucial role in regulatory
compliance and risk management, especially in setting minimum capital requirements under the
Basel Accords.

\textbf{The primary purpose of validating VaR models includes:}
\begin{enumerate}
\item Verifying if the VaR model provides reliable risk estimates.
\item Basel regulations mandate VaR validation to ensure sufficient capital is held by banks.
\end{enumerate}

\subsection{Key components of VaR}
a) Loss size \quad b) Probability \quad c) Time frame

\begin{mrfmlbox}
\[ \mathrm{VaR}(\alpha\%) \;=\; -\mu_{P/L} \;+\; \sigma_{P/L}\, z_\alpha \]
\mrvk{$\mu_{P/L}$, $\sigma_{P/L}$ = mean and standard deviation of the profit/loss
distribution;\quad $z_\alpha$ = the one-tailed standard normal deviate.}{%
\par\smallskip\textbf{\sffamily Not dynamic:} $\sigma_{P/L}$ is a single fixed number here, so
the measure carries no mechanism for volatility changing through time.}
\end{mrfmlbox}

\subsection{Challenges in traditional VaR models}
\begin{itemize}
\item It assumes past return distributions remain valid, ignoring dynamic volatility.
\item Many banks' historical VaR models have underperformed compared to dynamic models.
      Traditional models use historical returns but may not reflect changes in the portfolio.
\end{itemize}

To address this limitation, \term{GARCH VaR} models were developed to allow for more recent
volatility changes. The GARCH (Generalized Autoregressive Conditional Heteroskedasticity) model
helps address dynamic volatility by:
\begin{itemize}
\item adjusting for recent market conditions;
\item capturing changing risk levels in response to evolving market data.
\end{itemize}

Some studies indicate that the historical VaR models used by banks underperformed a basic model
based on GARCH VaR of the historical profit and loss (P\&L) of banks, which calls into question
the utility of historical VaR estimates.

\subsection{Pseudo history and pseudo returns in VaR modeling}
Traditional VaR models use historical portfolio return data, which ignores changes in portfolio
composition over time. This can misrepresent risk, especially when positions are actively
managed.

\begin{mrkeybox}[title={Why use pseudo returns?}]
\begin{itemize}
\item \term{Adjusts for portfolio changes} --- recalculates past returns based on the
      \textit{current} portfolio.
\item \term{More accurate risk estimation} --- captures real-time risk dynamics.
\item \term{Reflects risk reduction} --- if a trader hedges risk (e.g.\ with a derivatives
      overlay strategy), pseudo returns show lower VaR, while historical VaR remains unchanged.
\end{itemize}
\begin{center}
Pseudo returns $=$ more accurate VaR $=$ better risk management
\end{center}
\end{mrkeybox}

% ---------------------------------------------------------
\lo{MR-6 a}{Describe some important considerations for a bank in assessing the conceptual soundness of a VaR model during the validation process.}

\begin{mrexbox}[title={Example 1}]
A bank uses historical stock return data to estimate VaR, but it does not account for
derivatives hedging in its portfolio. The model then \term{overestimates} risk, leading to
excess capital requirements.
\end{mrexbox}

\begin{mrexbox}[title={Example 2}]
A trader selling deep out-of-the-money puts initially sees small, consistent profits, and a
P\&L-based VaR reflects low risk. However, if the underlying asset's price moves significantly
against the trader, the puts move closer to being in the money, and the potential loss becomes
substantial. VaR fails here as well.
\end{mrexbox}

\begin{mrdefbox}[title={Conceptual soundness}]
In both examples above, VaR has not provided proper guidance about the \term{real risk}
inherent. This is what we call conceptual soundness. VaR models must reflect changes in risk due
to position adjustments. If they fail to do so, they are conceptually unsound.
\end{mrdefbox}

\subsection{General validation criteria}
\begin{itemize}
\item \term{Assess methodology soundness} --- ensure the model uses a robust and appropriate
      approach.
\item \term{Evaluate input data quality} --- reliable and accurate data is crucial for valid
      risk estimation.
\item \term{Check model assumptions} --- assumptions should align with the portfolio's risk
      characteristics.
\item \term{Fit for purpose} --- the VaR model should align with the bank's specific risk
      management objectives (e.g.\ regulatory capital, internal risk assessment).
\end{itemize}

\subsection{Limitations of historical VaR models}
\begin{itemize}
\item Studies suggest that historical VaR models may underperform compared to alternative models
      like GARCH VaR.
\item Banks typically use historical portfolio returns in VaR models, but these returns do not
      adjust for portfolio composition changes over time. This limitation can lead to misleading
      risk estimates if the portfolio has changed significantly --- hence pseudo returns.
\end{itemize}

\subsection{Challenges with missing data}
Banks with large trading desks may face missing return data for certain assets. Some assets may
lack sufficient historical data or may not have clear value estimates. Missing data complicates
VaR estimation and can lead to inaccurate results.

\subsection{Quantitative tests for validation}
\begin{enumerate}
\item \term{Sensitivity analysis} --- measures the impact of data limitations and model
      assumptions.
\item \term{Confidence intervals} --- helps assess the accuracy of VaR estimates and potential
      errors.
\end{enumerate}

% ---------------------------------------------------------
\lo{MR-6 b}{Explain how to conduct sensitivity analysis for a VaR model, and describe the potential benefits and challenges of performing such an analysis.}

Sensitivity analysis for a VaR model examines how changes in inputs, particularly positions,
affect the model's output. This is vital for assessing the impact of simplifications or
omissions, especially in constructing the ``pseudo history'' of portfolio value changes. It
helps ensure these simplifications do not materially distort the VaR result, and offers a
structured way to prioritize model enhancements by incorporating more risk factors or
higher-order terms.

\begin{mrexbox}[title={Example of sensitivity analysis}]
A bank is concerned about the valuation of complex derivatives in its portfolio. Sensitivity
analysis can reveal how much the overall VaR would change if the valuation model for these
derivatives were adjusted --- say, by changing a key input parameter like volatility.
\end{mrexbox}

\subsection{1) Identify key inputs}
\begin{itemize}
\item Determine key inputs and assumptions (e.g.\ portfolio positions).
\item Validate assumptions and identify simplifications in pseudo-history.
\item Address omissions in pseudo-history, as they can materially affect VaR.
\item Track data proxies and monitor risks not included in the model. We use proxy data because
      sometimes asset data can be limited.
\end{itemize}

\subsection{2) Adjust key inputs and recalculate VaR}
\begin{itemize}
\item Adjust one input at a time (e.g.\ position weight).
\item Recalculate the VaR threshold.
\item \term{Marginal VaR:} the change in portfolio value due to a small change in a position's
      weight. Estimated as the slope coefficient in a regression of the change in the
      \textit{position's} value against the change in the \textit{portfolio's} value.
\end{itemize}

\begin{mrfmlbox}
\[ \Delta V_{it} \;=\; \alpha + \beta_i\, \Delta V_{pt} + \epsilon_t \]
\mrvk{$\Delta V_{it}$ = change in the value of position $i$ at time $t$ (the dependent
variable);\quad $\Delta V_{pt}$ = change in the value of the whole portfolio at time $t$ (the
independent variable);\quad $\beta_i$ = the slope coefficient, which is the marginal VaR
estimate;\quad $\alpha$ = intercept;\quad $\epsilon_t$ = residual.}{}
\end{mrfmlbox}

\subsection{3) Analyze results}
\begin{itemize}
\item Compare new VaR estimates to the original.
\item Evaluate VaR's sensitivity to changes in each input.
\item Sensitivity depends on:
  \begin{itemize}
  \item the change in the position's share of the portfolio;
  \item the sensitivity of the position's value to the overall portfolio value.
  \end{itemize}
\end{itemize}

\begin{center}
\small
\begin{tabularx}{\linewidth}{@{}p{3.4cm} >{\raggedright\arraybackslash}X@{}}
\arrayrulecolor{FRMRule}\toprule
\rowcolor{FRMNavy} \hdr{Benefits} & \hdr{What it looks like}\\
a) Risk assessment & Identifies the most influential factors in the portfolio's risk profile.
  \textit{Example:} identifying a specific trading desk's positions as highly sensitive to
  interest rate changes.\\
\rowcolor{FRMSoft}
b) Model validation & Provides a framework for testing the robustness of the VaR model against
  changes in assumptions, which enhances the model's credibility. \textit{Example:} testing the
  impact of different volatility assumptions on VaR.\\
c) Regulatory compliance & Helps meet regulatory requirements for model validation.\\
\rowcolor{FRMSoft}
d) Improved decision-making & Enables informed decisions regarding position adjustments and risk
  mitigation.\\
\midrule
\rowcolor{FRMNavy} \hdr{Challenges} & \hdr{What it looks like}\\
a) Data limitations & Requires sufficient data, which may be scarce. \textit{Example:} limited
  historical data for newly issued financial instruments.\\
\rowcolor{FRMSoft}
b) Computational complexity & Can be computationally intensive for large portfolios or complex
  models. \textit{Example:} Monte Carlo simulations with a large number of scenarios and risk
  factors.\\
c) Proxy issues & Using proxies can introduce inaccuracies. \textit{Example:} using a broad
  market index as a proxy for a specific stock with different risk characteristics.\\
\bottomrule
\end{tabularx}
\end{center}

% ---------------------------------------------------------
\lo{MR-6 c}{Describe the challenges a financial institution could face when calculating confidence intervals for VaR.}

When estimating a Value-at-Risk figure for a portfolio, assessing the accuracy of that estimate
is crucial. Statistically, this is done by placing confidence intervals around the VaR estimate,
indicating that the true VaR value should fall within this interval a certain percentage of the
time. This addresses the \term{estimation risk} inherent in VaR calculations.

While assessing estimation risk is standard in many contexts, many financial institutions do not
routinely calculate confidence intervals for their VaR estimates, despite VaR being a
statistical framework. This lack of confidence interval calculation is a notable gap in model
validation.

\subsection{1) Data quality and availability}
\begin{itemize}
\item \term{Historical data limitations.} Accurate confidence intervals require high-quality
      historical data. Incomplete, outdated, or erroneous data leads to unreliable VaR and
      confidence interval estimates.
\item \term{Non-normality of returns.} Financial returns often exhibit non-normal distributions
      (e.g.\ fat tails, skewness). This makes applying standard statistical methods, which
      assume normality, problematic, resulting in inaccurate confidence intervals.
\end{itemize}

\subsection{2) Model assumptions}
\begin{itemize}
\item \term{Choice of distribution.} Selecting the correct probability distribution is crucial.
      Using the wrong distribution leads to significant errors in confidence interval estimates.
\item \term{Alternatives to normality:}
  \begin{itemize}
  \item Base the confidence interval on the variance of the distribution of pseudo-historical
        returns, using more than a single point of the quantile estimate.
  \item Non-parametric approaches:
    \begin{enumerate}[label=\alph*)]
    \item \term{Order statistics} --- relies on the ordered history of pseudo-returns, just like
          the historical simulation method.
    \item \term{Bootstrap techniques} --- sampling with replacement from the pseudo-history to
          generate a bootstrapped pseudo-history, calculating VaR for each sample, and obtaining
          a distribution of VaR.
    \end{enumerate}
  \end{itemize}
\end{itemize}

\subsection{3) Historical findings}
\begin{itemize}
\item \term{Asymmetry.} Confidence intervals are not symmetric.
\item \term{Data quantity.} More data leads to tighter (more precise) confidence intervals.
\item \term{VaR methodology.} GARCH VaR of historical P\&L tends to produce \textit{tighter}
      confidence intervals compared to historical simulation VaR.
\item \term{Non-parametric methods.} Order statistics and bootstrap techniques do not
      necessarily produce tighter confidence intervals.
\item \term{Stress periods.} Historical simulation VaR produces \textit{wider} (less precise)
      confidence intervals during stress periods.
\end{itemize}

% ---------------------------------------------------------
\lo{MR-6 d}{Discuss the challenges in benchmarking VaR models and various approaches proposed to overcome them.}

Benchmarking, the comparison of a bank's VaR model against an alternative model, is often the
\term{most neglected} aspect of market risk model validation.

While parallel runs of new and old models are common for initial checks, formal comparisons are
rare due to the resource-intensive nature of building even a single VaR model. Constructing and
maintaining two models significantly increases the burden.

The most common benchmarking practice involves simply plotting the VaR outputs of two models
over time. This visual comparison provides limited insight, typically only revealing whether one
model is more conservative than the other or if they exhibit similar behavior. This lack of
formal statistical testing weakens the validation process.

\begin{mrtrapbox}[title={Issues with statistical tests that compare two VaR models}]
\begin{enumerate}
\item \term{Non-i.i.d.\ errors due to changing trading portfolios.} Trading portfolios change
      frequently, causing the errors (differences between predicted and realized P\&L) to
      violate the assumption of independence and identical distribution. This is particularly
      problematic for regression-based tests, which rely on the i.i.d.\ assumption.
\item \term{Lack of available alternative models.} A more practical challenge is the scarcity of
      readily available alternative VaR models for comparison. Banks rarely have two fully
      operational VaR models running simultaneously.
\end{enumerate}
\end{mrtrapbox}

\subsection{Proposed solutions}
\begin{enumerate}[label=\textbf{\alph*)}]
\item \term{Backtesting as benchmarking.} Since a second model is often unavailable, backtesting
      the existing VaR model against actual profit and loss outcomes can serve as a robust
      benchmarking tool. Specifically, compare the bank's positional VaR to the P\&L GARCH VaR.
\item \term{Positional VaR vs P\&L GARCH VaR:}
  \begin{itemize}
  \item \term{Positional VaR:} VaR based on current portfolio positions.
  \item \term{P\&L GARCH VaR:} VaR calculated using a GARCH model applied to the bank's
        historical P\&L.
  \end{itemize}
\end{enumerate}

\begin{mrkeybox}[title={Empirical findings of backtesting positional VaR vs P\&L GARCH VaR}]
\begin{enumerate}
\item \term{Conservatism.} Positional VaR tends to be \textit{more conservative} than P\&L GARCH
      VaR due to its link to regulatory capital requirements.
\item \term{Accuracy.} Positional VaR \textit{rarely outperforms} P\&L GARCH VaR in terms of
      accuracy, likely because the former is more conservative.
\end{enumerate}
\end{mrkeybox}


% ==========================================================
